\subsection{What is not established here}\label{sec:doubling_differences}

Three claims that would naturally accompany the results above are absent, and in each case for
the same reason: they are properties of the backward chain as a stochastic process, and no such
process is constructed.

\begin{itemize}
  \item \emph{Transience and null recurrence.} The phase diagram distinguishes only between the
        existence and the non-existence of an invariant probability. Where none exists --- at
        $s<1$, and at $s=1$ with $c\geq1$ --- nothing is said about which of transience or null
        recurrence obtains. In particular the case $s=1$, $c=1/\ln2$ sits inside the
        non-existence regime and is not resolved further.
  \item \emph{An infinite expected hitting time.} At $s=1$ with $1\leq c<2$ the absence of an
        invariant probability is established; the accompanying statement that
        $\mathbb E(\sigma\mid X_0=j)=+\infty$ at every ladder state is not.
  \item \emph{The trajectory reading of the diffusion operator.} The operator excluded at $p=2$
        is characterised by the resolvent identities, not as the sum
        $\sum_{n\geq0}\bigl(\mathbb E(\theta(X_n)\mid X_0=x)-\lambda(\theta)\bigr)$ along
        backward trajectories. That the two agree needs the chain.
\end{itemize}

\noindent
Nothing above depends on these. They are named so that a reader who expects the vocabulary of
recurrence classification knows where it has been replaced by the weaker statement that is
actually proved.
